%% cattree.tex -- "The cat gym": the theme=cattree margin describes itself. LuaLaTeX:
%%   lualatex cattree.tex
\documentclass[11pt,a4paper]{article}
\renewcommand\rmdefault{cmr}\renewcommand\sfdefault{cmss}\renewcommand\ttdefault{cmtt}
\usepackage[T1]{fontenc}
\usepackage[a4paper,left=68mm,right=22mm,top=32mm,bottom=30mm,headsep=8mm]{geometry}
\usepackage{fancyhdr}
\usepackage[grown,sections,theme=cattree]{florilegium}

\pagestyle{fancy}\fancyhf{}
\fancyhead[R]{\small\textsc{The cat gym}}
\fancyfoot[C]{\small\thepage}
\renewcommand\headrulewidth{0pt}
\newcommand\code[1]{\texttt{\small #1}}

\begin{document}
\section{A gym for every page}

\florilegiuminitial[4]{T}{he} margin of this page is a cat gym, and the cat on the letter
that opens this paragraph is the first of its residents. The package draws it with
\code{theme=cattree}: the same ink as the flowering plant, washes of one blue from navy to
mist with copper for the fine work, but the plant is replaced by posts wound with sisal,
plush platforms, condos, hexagons, hammocks and rope bridges.

Nothing is stored. Every piece is computed while the page is written, at one twentieth of
its real size: a post is nine centimeters thick, a hexagon sixty across, a condo's door
twenty. A gray cat of forty-five centimeters was the yardstick while the pieces were drawn,
so that every hexagon and basket really holds a cat.

\section{Heavy below, light above}

The gym is built the way a cat would climb it. On the floor stands a tower of three to five
posts, in whole levels about fifty centimeters apart. Where two neighbors meet, the shorter
one holds up a platform and the taller one runs on through it. Crossbeams tie the posts
together low down, hammocks hang under the high platforms, and condos, baskets on arms,
rings and sisal tubes take what is left.

Above the tower the gym goes onto the wall: a route of boards, a rope bridge to a half-moon
shelf, never steeper than fifteen degrees and only ever between two raised points, a
hexagon or two and a clear bubble. Near the top it thins out to a few half-moons, joined by
floating steps, pegs and a line of dots growing smaller as they rise.

\section{Filling the space}

Then the program does what the garden does. It lays a grid over the margin, finds the
largest empty space and puts something there, a hexagon or a half-moon, joined to the gym
the way a cat would get there: floating steps when it is level, dots, pegs or a gangway on a
slant, and going up a ladder hung flat on the wall, a pole with half-round steps, or nothing,
the shelves staggered for a jump. Every piece keeps a few centimeters from every
other, so nothing passes through anything. Each heading gets a cat on a shelf beside it,
meowing.

Last come the cats, one or two to a page, so that most of the gym stands ready for them.
They settle where they fit: lying on a platform, curled in a hammock or a basket, sitting in a
hexagon, stretching, scratching a post, draped over an edge, peeking out of a condo, walking
across a bridge.

\section{The last page}

Where the text stops early, as it does here, the gym spreads along the floor under it, an
L on posts of its own, and often a U: a second tall post at the far end with a bridge back
to the tower.
\florilegiumtextend
\end{document}
